\documentclass[10pt,a4paper]{amsart}
\usepackage[margin=19mm]{geometry}
\usepackage{amsmath,amssymb,mathtools}
\usepackage[scr=rsfs,cal=euler]{mathalfa}
\usepackage{xfrac}
\usepackage[unicode,hidelinks,bookmarks=false]{hyperref}
\newcommand{\ie}{i.e.\ }
\newcommand{\cf}{cf.\ }
\newcommand{\resp}{respectively\ }
\newcommand{\Subsection}{\subsection}
\providecommand{\nobreakdash}{\nobreak\hbox{-}}
\setlength{\emergencystretch}{2em}
\multlinegap0pt
\usepackage{bm}
\usepackage{bbm}
\usepackage{dsfont}
\usepackage{xspace}
\usepackage{cancel}
\usepackage{mathtools}
\usepackage{amsthm}
\makeatletter
\@ifundefined{theorem}{\newtheorem{theorem}{Theorem}}{}
\makeatother
\makeatletter
\expandafter\ifx\csname keywords\endcsname\relax
\newenvironment{keywords}
  {\par\noindent\textbf{Keywords: }} %\itshape}
  {\par}
\fi
\makeatother
\title[The Limits of AI Explainability: An Algorithmic Information ]{The Limits of AI Explainability: An Algorithmic Information Theory Approach}
\author{Shrisha Rao}
\date{Source: arXiv:2504.20676v2 (CC BY 4.0)}
\begin{document}
\maketitle

\noindent\emph{Source and licence.} The Limits of AI Explainability: An Algorithmic Information Theory Approach. Version arXiv:2504.20676v2 (CC BY 4.0), distributed under the Creative Commons Attribution 4.0 International licence, as recorded in the arXiv licence metadata for this exact version and shown on the abstract page https://arxiv.org/abs/2504.20676v2. Full rights notice: licenses/arxiv-2504.20676-CC-BY-4.0.txt.\par\smallskip

\noindent\textit{Editorial scope.} Counted unit: the Theorem whose source label is \texttt{thm:local-complexity} in the article (source0.tex lines 1278--1289, source label \texttt{thm:local-complexity}), delivered together with the article's own complete original proof of it (source0.tex lines 1291--1371, one \texttt{proof} environment of 81 lines). No mathematics is written, repaired or re-derived: statement and proof are the article's own text line for line. Required context retained uncounted: none. Every definition, display and label of those lines is retained unchanged. Omitted from this package and not redistributed: 1--1277, 1290--1290, 1372--2410. Nothing inside a retained excerpt is omitted or altered: every retained body is delivered exactly as the article's lines are written, so no omission receipt of any kind accompanies this package. Citations: the retained excerpts carry no citation command at all, so no bibliography entry of the article is transcribed and no reference is redistributed. Reference adaptation: none was needed and none was made; every reference inside the delivered unit resolves inside it, so no unresolved reference or citation is produced. Printed slips kept literal: no printed word, symbol, hypothesis or number of the counted statement or of its proof was altered. Layout and class: the article's own class is \texttt{article} with packages including amsmath, amssymb, amsthm, algorithm, algpseudocode, graphicx, hyperref, xcolor, mathtools; the wrapper is the lane's pinned \texttt{amsart} editorial wrapper at 10pt on A4, which restates the theorem-like environments the retained text needs and adds the article's own macro definitions that the retained text needs (none), with extra wrapper packages (bm, bbm, dsfont, xspace, cancel, mathtools). The delivered numbering is the unit's own. Assets: no retained span contains a figure environment, a graphics-inclusion command, a PDF-page inclusion, a table or a TikZ picture; no image file is copied into this package, the package declares no asset dependency and no image file is redistributed with it. Licence: version arXiv:2504.20676v2 of this article is distributed under the Creative Commons Attribution 4.0 International licence, as recorded in the arXiv licence metadata for this exact version and shown on the abstract page https://arxiv.org/abs/2504.20676v2. A full rights notice is delivered at licenses/arxiv-2504.20676-CC-BY-4.0.txt. Characters: every retained excerpt is pure ASCII, so the delivered engine, XeLaTeX with its default Latin Modern text font, reports no missing character.
\begin{theorem}[Local Explanation Complexity]
\label{thm:local-complexity}
Let $\mathcal{X} \subseteq \mathbb{R}^d$ be a compact convex set, let $f$ be an $L$-Lipschitz continuous function on $\mathcal{X}$, and let $N(x_0) \subseteq \mathcal{X}$ be a neighborhood of $x_0 \in \mathrm{int}(\mathcal{X})$. Let $r > 0$ be such that $B_r(x_0) \supseteq N(x_0)$ and $B_r(x_0) \subseteq \mathcal{X}$. Then the local complexity function satisfies:
\begin{align}
\kappa_f^{\text{local}}(\delta, x_0, N) = 
\begin{cases}
\mathcal{O}(1) & \text{if } \delta \geq Lr \\
\mathcal{O}(d \log(Lr/\delta)) & \text{if } \delta < Lr
\end{cases}
\end{align}
where $\kappa_f^{\text{local}}(\delta, x_0, N)$ is the minimum description length of a program that, given oracle access to $f$, computes a local approximation $g$ with error at most $\delta$ in the neighborhood $N(x_0)$.
\end{theorem}

\begin{proof}
Since $N(x_0) \subseteq B_r(x_0) \subseteq \mathcal{X}$ and $f$ is $L$-Lipschitz on $\mathcal{X}$, the function $f$ is also $L$-Lipschitz on $B_r(x_0)$. We analyze the local explanation complexity by considering two cases based on the relationship between the desired accuracy $\delta$ and the maximum variation $Lr$ of the function within the ball.

\textbf{Case 1: $\delta \geq Lr$.}

By the Lipschitz property, for any $x \in B_r(x_0) \subseteq \mathcal{X}$, we have:
\begin{align}
|f(x) - f(x_0)| \leq L \|x - x_0\|_2 \leq Lr
\end{align}

Therefore, the constant function $g(x) = f(x_0)$ achieves:
\begin{align}
\sup_{x \in N(x_0)} |f(x) - g(x)| \leq \sup_{x \in B_r(x_0)} |f(x) - g(x)| \leq Lr \leq \delta
\end{align}

To specify this approximation, we require a program that takes input $x \in N(x_0)$ and returns $f(x_0)$. Since we measure relative complexity (description length given oracle access to $f$), the program need only specify:
\begin{itemize}
\item The algorithm: ``Compute and return $f(x_0)$ for all inputs in $N(x_0)$''
\item The point $x_0 \in \mathcal{X}$, which requires $\mathcal{O}(\log(1/\epsilon))$ bits for $\epsilon$-precision encoding
\item The constant value $f(x_0)$ is obtained via oracle access
\item The total encoding requires $\mathcal{O}(1)$ bits (treating the precision as a constant)
\end{itemize}

Thus, $\kappa_f^{\text{local}}(\delta, x_0, N) = \mathcal{O}(1)$ when $\delta \geq Lr$.

\textbf{Case 2: $\delta < Lr$.}

Within $B_r(x_0)$, the function $f$ may vary by up to $Lr > \delta$, so we require a more refined approximation. We construct a piecewise constant approximation using a grid-based approach on $B_r(x_0)$.

\textbf{Grid construction:} By Lipschitz continuity, any region of diameter $d'$ has variation at most $Ld'$. To ensure variation at most $\delta$, we need:
\begin{align}
Ld' \leq \delta \quad \Rightarrow \quad d' \leq \frac{\delta}{L}
\end{align}

Cover $B_r(x_0)$ with balls of radius $\rho = \delta/(2L)$. The number of such balls required to cover $B_r(x_0)$ is at most:
\begin{align}
N_{\text{balls}} = \mathcal{O}\left(\left(\frac{r}{\rho}\right)^d\right) = \mathcal{O}\left(\left(\frac{Lr}{\delta}\right)^d\right)
\end{align}

This follows from standard covering number bounds for balls in $\mathbb{R}^d$.

\textbf{Approximation construction:} For each small ball $B_i$ in the cover, choose a representative point $x_i \in B_i \cap B_r(x_0)$ and define:
\begin{align}
g(x) = f(x_i) \quad \text{for } x \in B_i \cap B_r(x_0)
\end{align}

For any $x \in B_i \cap B_r(x_0)$, we have $\|x - x_i\|_2 \leq 2\rho = \delta/L$, so:
\begin{align}
|f(x) - g(x)| = |f(x) - f(x_i)| \leq L\|x - x_i\|_2 \leq L \cdot \frac{\delta}{L} = \delta
\end{align}

Since $N(x_0) \subseteq B_r(x_0)$, this bound holds for all $x \in N(x_0)$.

\textbf{Complexity analysis:} To encode the piecewise constant function $g$, given oracle access to $f$, we need:
\begin{enumerate}
\item \textbf{Grid specification:} Encode the partition of $B_r(x_0)$ into regions. This requires specifying the covering scheme and resolution $\rho$, which takes $\mathcal{O}(\log(Lr/\delta))$ bits.

\item \textbf{Representative points:} For each of the $\mathcal{O}((Lr/\delta)^d)$ regions, we store a representative point $x_i$. Each point requires $\mathcal{O}(d \log(Lr/\delta))$ bits for sufficient precision.

\item \textbf{Function values:} The values $f(x_i)$ are obtained via oracle queries and need not be explicitly encoded in the program describing $g$.

\item \textbf{Lookup algorithm:} A program that, given $x \in N(x_0)$, determines which region contains $x$ and returns the corresponding constant value. This requires $\mathcal{O}(1)$ bits for the algorithm description.
\end{enumerate}

The total description length is:
\begin{align}
K(g) &= \mathcal{O}(\log(Lr/\delta)) + \mathcal{O}\left(\left(\frac{Lr}{\delta}\right)^d \cdot d \log(Lr/\delta)\right) \\
&= \mathcal{O}\left(\left(\frac{Lr}{\delta}\right)^d d \log(Lr/\delta)\right)
\end{align}

However, we can improve this bound by noting that for local explanations, we do not need to encode the global partition, only the local behavior. Using a more efficient encoding scheme based on recursive subdivision (e.g., quadtree-like structures adapted to $B_r(x_0)$), we can achieve:
\begin{align}
K(g) = \mathcal{O}(d \log(Lr/\delta))
\end{align}

This follows from the fact that the depth of recursion is $\mathcal{O}(\log(Lr/\delta))$ and at each level we need $\mathcal{O}(d)$ bits to specify the subdivision.

Thus, $\kappa_f^{\text{local}}(\delta, x_0, N) = \mathcal{O}(d \log(Lr/\delta))$ when $\delta < Lr$.

\textbf{Comparison with global complexity:} For global explanations on $\mathcal{X}$, if $\mathcal{X}$ has diameter $D$, the complexity scales as $\mathcal{O}((LD/\delta)^d d \log(LD/\delta))$, which is exponential in $d$. In contrast, the local complexity is only logarithmic in the ratio $Lr/\delta$, demonstrating that local explanations can be exponentially simpler than global ones.
\end{proof}

\end{document}
